\section{Our Proposed Scheme}
\label{sec:scheme}

The scheme consists of an offline phase, run once before training, and an online phase of training
rounds. This section describes the offline phase. Its first step builds the \emph{union list}: one
entry per distinct concept held by any client, such that labels of different clients that denote the
same concept share one entry. The list is built without revealing any client's label set
$\mathcal{Y}_i$ to the other clients or to the aggregator.

\subsection{Union List by Circuit PSI}
\label{sec:union}

\paragraph{Goal}
For every pair of labels $y_i\in\mathcal{Y}_i$ and $y_j\in\mathcal{Y}_j$ of different clients, the
clients must decide whether $y_i$ and $y_j$ denote the same concept, i.e., whether both
$d_{\mathrm{sem}}(y_i,y_j)$ and $d_{\mathrm{vis}}(D_i^{(y_i)},D_j^{(y_j)})$ are small. A plain PSI
would reveal which labels match. We therefore use circuit PSI: the pairwise matches and the
resulting groups stay inside an MPC among the clients, and the only output is, for every label, a
random group secret $K$ that is equal for all labels of one concept and is revealed only to the
label's owner. From these secrets the union list is assembled by a secure sum
(Sec.~\ref{sec:union-secsum}).

\paragraph{Semantic description of a label}
Each client $C_i$ describes every label $y\in\mathcal{Y}_i$ by a keyword $w_i(y)$, a short word or
phrase naming the concept of $D_i^{(y)}$, chosen locally. Two matching modes are supported.
\begin{itemize}
  \item \emph{Exact mode.} When the application domain shares a public vocabulary (for example, all
  clients are hospitals using a common medical terminology), the clients take $w_i(y)$ from a
  \emph{public dictionary} $\mathcal{W}$ of canonical terms, with synonyms folded into one entry.
  Two labels then match semantically iff their keywords are equal, which the MPC tests on hashed
  values $H_{\mathrm{n}}(w_i(y))$.
  \item \emph{Similarity mode.} When the clients do not agree on a common vocabulary, or do not want
  to commit their categories to one, each client writes its keyword freely. A public text encoder
  $E$ maps a keyword to a unit vector $\mathbf{e}=E(w_i(y))$, and the semantic test is the CSLS
  criterion
  \begin{equation}
    2\langle\mathbf{e},\mathbf{e}'\rangle-r-r'\ \ge\ \tau ,
    \label{eq:csls}
  \end{equation}
  where $r$ is the mean similarity of $\mathbf{e}$ to its $h$ nearest words of a public anchor
  vocabulary. The hub terms $r,r'$ discount generic words that are close to everything, and the
  threshold $\tau$ is public.
\end{itemize}

\paragraph{Visual description of a label}
A keyword alone cannot tell, e.g., a handwritten ``3'' from a printed one, or a sketch of a cat from
a photograph. Each label is therefore also described by its samples, without revealing them. From a
public seed, every client generates the same \emph{probe set}: $N_I$ random images
$\boldsymbol\xi_1,\dots,\boldsymbol\xi_{N_I}$, together with their coordinates under a fixed public
signature function $\sigma$ (e.g., a randomly initialized convolutional feature map). For a label
$y$, client $C_i$ computes the mean signature $\bar{\boldsymbol\sigma}_i(y)$ of the samples in
$D_i^{(y)}$ and records which $k$ probes lie nearest to it:
\begin{equation}
  \mathbf{s}_i(y)\in\{0,1\}^{N_I},\qquad \mathbf{s}_i(y)[j]=1 \iff
  \boldsymbol\xi_j \text{ is among the $k$ nearest probes to } \bar{\boldsymbol\sigma}_i(y).
\end{equation}
Two labels are visually similar if their probe sets overlap,
$\langle\mathbf{s}_i(y_i),\mathbf{s}_j(y_j)\rangle\ge t$, a coarse and public test that approximates
a small $d_{\mathrm{vis}}$. Only the $k$ probe indices enter the computation; neither samples nor
signatures leave the client.

\paragraph{Secure matching and grouping}
Each client turns each label into a row $(v,\phi,\mathbf{s})$ with a valid bit $v$, the keyword
feature $\phi$ (hashed keyword or fixed-point embedding with its hub term), and the probe
indicator $\mathbf{s}$; dummy rows with $v=0$ pad every client to the same number of rows $m$, so the
row count hides $|\mathcal{Y}_i|$. The clients Shamir-share their rows among themselves and evaluate,
inside an honest-majority MPC, the match bit of every pair of rows $a,a'$ of different clients,
\begin{equation}
  \mu_{aa'} = v_a v_{a'}\,\big[\langle\mathbf{s}_a,\mathbf{s}_{a'}\rangle\ge t\big]\,
  \mathrm{KW}(\phi_a,\phi_{a'}),
\end{equation}
where $\mathrm{KW}$ is keyword equality (exact mode) or test~\eqref{eq:csls}. Matching
is not transitive, so the MPC then merges matching rows into connected components by min-label
propagation; every row carries a joint random secret, and every row of a component ends with the
secret of its root. Finally each owner receives the secret $K$ of its own rows and nothing else.
Labels of one concept, at any clients, now hold the same $K$; no party sees a match bit, a group, or
another client's secrets, and since $K$ is random, it says nothing about the label it stands for.

\subsection{Delivering the Union List to the Aggregator}
\label{sec:union-secsum}

The aggregator should learn the union list but not who holds which entry. Each client maps every
secret $K$ it holds to a bucket $\pi(K)=H_{\mathrm b}(K)\in[2^\beta]$ and builds a vector
$\mathbf{y}_i$ of $2^\beta$ entries: a uniformly random nonzero value $\alpha$ in each of its buckets,
together with a check value $\alpha\cdot\mathrm{chk}(K)$, and zeros elsewhere. The clients run the
secure sum of Sec.~\ref{sec:prelim} (secure aggregation), $\boldsymbol\Sigma=\mathsf{SecSum}(\{\mathbf{y}_i\})$,
which gives the aggregator $\sum_i\mathbf{y}_i$ and nothing else. Holders of one concept write into
the same bucket, so the nonzero entries of $\boldsymbol\Sigma$ are exactly the union list; since a sum
of random values is random, it hides how many clients hold each entry. The aggregator posts the
occupied buckets, sorted, on $\mathsf{BB}$; their number $L$ is the size of the union. A bucket that
received two different secrets fails the check $\boldsymbol\Sigma[\pi]_1=\mathrm{chk}(K)\,
\boldsymbol\Sigma[\pi]_0$ at its holders, and the step is repeated with a fresh public salt. Because
buckets are hashes of secrets known only to their holders, the posted list reveals no label to the
aggregator or to any other client; each client finds the position of its own buckets in the list,
which becomes the index of the corresponding label in the next step.
